
En el proceso de producción musical, la mezcla de audio representa una de las etapas más complejas y determinantes para la calidad final de una obra sonora. Dentro de ese proceso, la mezcla vocal ocupa un lugar central en la mayoría de los géneros musicales populares, pues, la voz es el elemento principal y su tratamiento técnico a través de cadenas de ecualización, compresión, de-essing, saturación y efectos espaciales influye directamente en la claridad, presencia e integración del canto dentro del instrumental. Sin embargo, construir una cadena de procesamiento vocal coherente y efectiva exige conocimientos técnicos que van más allá del uso básico de un DAW. Como señalan De Man, Reiss y Stables en [4], incluso con acceso a herramientas digitales de alta calidad, un productor sin experiencia formal casi inevitablemente generará problemas sonoros durante la grabación y la mezcla, lo que aumenta la necesidad de sistemas que puedan orientar el proceso.

Esta situación se ha vuelto especialmente relevante en el contexto de la producción musical independiente. Con el auge de las herramientas de producción como softwares, plugins y DAWs accesibles a cualquier persona con un computador, una cantidad creciente de productores musicales asumen simultáneamente los roles de grabación, edición y mezcla sin contar con una formación técnica sólida en cada una de esas áreas. Vanka et al. en [1] documentaron esta realidad a través de entrevistas y cuestionarios con ingenieros de distintos niveles de experiencia y confirmaron que los productores amateurs perciben la mezcla como el principal obstáculo para publicar su música, y que muchos trabajan mediante ensayo y error al no contar con una guía técnica clara para organizar sus cadenas de procesamiento.

La literatura sobre sistemas de mezcla automática e inteligente ha explorado diversos caminos para abordar este problema. Por un lado, Martínez Ramírez y Reiss en [5] demostraron que las redes neuronales profundas pueden aprender transformaciones de audio complejas directamente desde señales crudas, sin necesidad de modelar efectos individuales de manera explícita. Por otro lado, Martínez Ramírez, Stoller y Moffat en [2] confirmaron que este tipo de sistemas puede producir mezclas perceptualmente indistinguibles de las realizadas por ingenieros humanos, al menos en contextos instrumentales bien delimitados como la batería. Sin embargo, estos sistemas están orientados principalmente a la automatización completa del proceso de mezcla, lo cual no se ajusta necesariamente a las necesidades de todos los usuarios ni respeta la naturaleza creativa y subjetiva del proceso.

En ese marco, Moliner et al. en [3] plantearon una crítica directa a los enfoques determinísticos: tratar la mezcla como un problema de regresión de uno a uno ignora el hecho de que múltiples soluciones válidas pueden existir para las mismas pistas de entrada, lo que resulta en sistemas que producen mezclas planas y que no tienen en cuenta la diversidad de decisiones artísticas posibles. Este planteamiento refuerza la idea de que la mezcla vocal, con su alta variabilidad técnica y estilística y su dependencia de las características acústicas específicas de cada voz, no puede ser tratada como un problema de una única solución técnica universal.

Por su parte, Vanka et al. en [1] aportaron evidencia empírica de que los distintos perfiles de usuario requieren soluciones diferentes: mientras los amateurs prefieren sistemas completamente autónomos, los semi-amateurs y los profesionales valoran más las herramientas que funcionan como asistentes, ofreciendo un punto de partida que pueden ajustar y refinar según su criterio. Esta distinción es fundamental para el diseño del sistema propuesto en el presente proyecto.

Como consecuencia de lo anterior, existe una necesidad concreta de desarrollar herramientas que no busquen reemplazar el criterio del ingeniero de sonido, sino que lo asistan de manera técnicamente informada. Una plataforma capaz de analizar las características acústicas de una señal vocal y proponer cadenas iniciales de procesamiento ajustables respondería directamente a esa necesidad: reduciría la incertidumbre técnica que enfrentan los productores con menor experiencia y facilitaría el aprendizaje del proceso de mezcla y permitiría al usuario mantener el control creativo sobre las decisiones finales.
